\newpage

\vspace*{1cm}

\begin{center}
    \textbf{Abstract}
\end{center}

\vspace*{1cm}

\noindent 

Diese Arbeit befasst sich mit der Weiterentwicklung und Evaluation eines am KI-Campus eingesetzten Retrieval-Augmented-Generation-Chatbots im Bildungskontext. Ziel ist es, die bestehende Systemarchitektur und die zugrunde liegende Datenverarbeitung so zu verbessern, dass kontextgestützte Antworten innerhalb eines Lernmanagementsystems qualitativ gesteigert werden.

Grundlage der Weiterentwicklung ist die Analyse von rund 13.000 realen Nutzeranfragen, anhand derer zentrale Nutzungsmuster identifiziert wurden. Daraus werden drei Use Cases abgeleitet: technischer Support, kursbezogene Fachfragen und lernunterstützende Dialoge. Aufbauend auf diesen Erkenntnissen wird die Systemarchitektur gezielt erweitert. Die Implementierung umfasst den Ausbau der Wissensbasis, die Integration multimodaler Inhalte, eine routingbasierte Retrieval-Logik mit Single- und Multi-Hop-Verfahren sowie den Einsatz von Reranking-Mechanismen. Ergänzend werden die Prompt- und Zitierlogik überarbeitet, um Transparenz und Nachvollziehbarkeit der Antworten zu erhöhen.

Die Evaluation der weiterentwickelten Systemversion erfolgt mehrdimensional. Neben automatisierten Qualitätsmetriken werden ein expertenbasierter Blindvergleich im A/B-Design sowie Analysen von Latenzzeiten, Tokenverbrauch und Kosten durchgeführt. Die Ergebnisse zeigen eine verbesserte Kontexttreue, höhere Antwortrelevanz und eine gesteigerte didaktische Qualität. Gleichzeitig wird deutlich, dass die erhöhte Systemkomplexität mit einem höheren Ressourcenverbrauch verbunden ist.

Insgesamt verdeutlicht die Arbeit, dass die Kombination aus technischer Optimierung, empirischer Evaluation und ökonomischer Betrachtung eine fundierte Grundlage für die Weiterentwicklung RAG-basierter Chatbots im Bildungskontext schafft.